\chapter{大语言模型}
\label{chap:large-language-models}

\setcounter{footnote}{0}

\newcommand{\chTwentyBox}[1]{%
  \fbox{\rule{0pt}{1.65em}\makebox[7.2em][c]{\small #1}}%
}
\newcommand{\chTwentySmallBox}[1]{%
  \fbox{\rule{0pt}{1.35em}\makebox[4.8em][c]{\scriptsize #1}}%
}
\newcommand{\chTwentyCell}[1]{%
  \fbox{\rule{0pt}{1.25em}\makebox[1.7em][c]{\scriptsize $#1$}}%
}
\newcommand{\chTwentyTextCell}[1]{%
  \fbox{\rule{0pt}{1.25em}\makebox[1.7em][c]{\scriptsize #1}}%
}
\newcommand{\chTwentyNewCell}[1]{%
  \fcolorbox{black}{gray!25}{%
    \rule{0pt}{1.25em}\makebox[1.7em][c]{\scriptsize $#1$}}%
}
\newcommand{\chTwentyArrow}{\ensuremath{\longrightarrow}}

在 Google 研究人员于 2017 年提出 Transformer \cite{ch20-1} 之前，语言任务主要建立
在 IBM 开创的统计语言建模技术之上（例如 \cite{ch20-2}）。2018 年，Google 推出了
基于 Transformer 的 BERT 语言模型，并在自然语言任务中取得了前所未有的结果。此后，
Transformer 架构\footnote{本章中的“架构”指神经网络架构，而不是本书其他章节中所说的
硬件架构。}成为基于神经网络的自然语言处理（natural language processing，NLP）深度
学习方法的基础。

此后，大语言模型（large language models，LLM）受到产业界广泛关注，并迅速发展，在
文本生成、摘要、翻译和对话等多种任务中表现出接近人类的能力。与 BERT 及更早一代的
语言模型相比，这些模型的规模增长了数个数量级，以便容纳更多训练数据。“大语言模型”
这一术语正是为了把它们与早期语言模型区分开来。LLM 最初主要用于处理和生成文本信息，
但后来也出现了多模态模型：这类模型使用视觉、音频以及科学数据等不同类型的数据进行
训练，并在这些数据上工作。

LLM 的核心思想是：给定一串此前出现的 token，预测下一个 token。这里的 token 可以是
一个单词、子词或字符。在 OpenAI ChatGPT \cite{ch20-3} 这样的聊天机器人应用中，此前
的 token 序列统称为上下文（context），它由系统提示词、用户输入以及模型已经生成的
token 组成。为了生成正确、有效或有用的结果\footnote{对于文本摘要等任务，“准确率”
未必是合适的评价指标，因为不同答案可能都正确。本章使用“准确”时，采用更一般的
“正确、有效或有用”的含义。}，LLM 通常包含数十亿个参数，例如数十层、每层数亿个
参数，并且需要在规模巨大的数据集上训练。在推理阶段，更长的上下文有助于生成更准确
的响应。模型需要通过一种称为注意力（attention）的机制维护上下文，并识别大段上下文
中与下一个 token 生成最相关的部分。因此，理解注意力机制以及 Transformer 架构，对于
理解 LLM 十分重要。Transformer 是目前把注意力机制引入 LLM 的最主流方法之一
\cite{ch20-1}。

\noindent\textbf{补充说明：}虽然本章聚焦 Transformer 和注意力机制，研究界也在持续
探索具有不同优势的其他 LLM 架构，例如状态空间模型 Mamba \cite{ch20-4}、卷积模型
Hyena \cite{ch20-5} 以及混合模型。它们在某些场景下可能具有更低的内存需求。

\section{Transformer 架构}
\label{sec:ch20-transformer}

如今大多数 LLM 都采用 Transformer 架构及其注意力机制，以捕获 token 之间的依赖关系，
并判断这些 token 在生成下一个 token 时的重要程度。注意力机制类似于神经网络内部的
分类过程：它使用 softmax 函数生成概率，而序列中每个 token 获得的概率可以看作该 token
对于生成下一个 token 的相对重要性。

Transformer 最初是为机器翻译等序列转换任务提出的，因此包含两种类型的层。编码器层
对输入 token 进行上下文化（或混合），把源语言中的符号序列（例如 token）转换成连续
表示，即浮点嵌入向量序列。解码器层接收编码器输出的连续表示，并逐个生成目标语言中的
输出 token。解码器不仅接收编码器的输出，还接收自身最近生成的 token。因此，Transformer
模型是自回归（auto-regressive）模型。

Transformer 提出之后，人们把它适配到了其他任务。根据具体任务，一个 Transformer
不一定需要同时包含两类层。例如，执行分类等判别任务的模型并不需要解码器，因为它们
不会生成目标语言中的文本。另一方面，仅解码器模型（decoder-only model）能够执行文本
生成等生成式任务。LLM 就属于这一类：它根据用户查询生成文本响应，而不是执行序列到
序列的转换。生成式任务需要多次执行模型；每次执行都根据已有 token 序列，预测候选
新输出 token 的可能性。本节重点讨论仅解码器 Transformer 的架构。关于编码器--解码器
Transformer，可参考文献 \cite{ch20-1}。

\begin{figure}[htbp]
  \centering
  \small
  \begin{tabular}{c}
    \chTwentyBox{Tokenizer（分词器）}
    \quad $\downarrow$ \quad
    \chTwentyBox{Embedding（嵌入）}
    \quad $\downarrow$ \quad
    \chTwentyBox{位置编码}\\[0.5em]
    $\downarrow$\\[-0.2em]
    \fbox{\begin{minipage}{0.75\textwidth}
      \centering
      \textbf{Transformer 层（重复多次）}\\[0.3em]
      \chTwentySmallBox{多头注意力}
      \quad $\longrightarrow$ \quad
      \chTwentySmallBox{Add \& Norm}
      \quad $\longrightarrow$ \quad
      \chTwentySmallBox{前馈网络}
      \quad $\longrightarrow$ \quad
      \chTwentySmallBox{Add \& Norm}\\[0.35em]
      残差连接与层归一化贯穿各子层
    \end{minipage}}\\[0.5em]
    $\downarrow$\\[-0.2em]
    \chTwentyBox{Un-embedding（反嵌入）}
    \quad $\longrightarrow$ \quad
    \chTwentyBox{下一个 token 的概率分布}
  \end{tabular}
  \caption{仅解码器 Transformer 的基线架构（依据原书图 20.1 临时重绘，未包含
  KV cache）。}
  \label{fig:20-1}
\end{figure}

图 \ref{fig:20-1} 展示了一个基于仅解码器 Transformer 的 LLM 设计。下面简要说明
其中的主要组成部分。

\begin{enumerate}[label=\arabic*.,leftmargin=2.7em]
  \item \textbf{Tokenizer（分词器）}把输入文档或查询中的文本转换为 token 序列，即
  用来表示一个字符或一小段字符的整数 ID。文本到 token 的映射建立在这样的思想上：
  可以从目标语言的训练文档中学习出一个有限词表。分词器学习并识别出的所有不重复
  token 构成 LLM 的词表。每个分词器都有自己对 token 的定义，以及把文本输入转换成
  token 的方法。

  在训练 LLM 时，模型会为训练文档中遇到的 token 分配唯一 ID。每当分词器决定从下
  一个文本片段形成一个 token 时，它都会检查相同片段是否已经出现过。如果没有，就发现
  了一个新的文本片段，并为它分配新的 ID；这个 ID 将作为该片段在模型中的表示。文本
  片段与 token ID 的映射会被记录下来，供以后查找。如果片段此前已经出现，分词器就
  通过映射找到原先分配的 token ID。文本片段到 token ID 的映射是训练过程的结果之一。

  在推理阶段，每次分词器从文本片段形成 token 时，都会查找该片段对应的 token ID，
  并把它作为分词器的下一个输出。严格地说，token 是构成 LLM 词表的文本片段，而
  token ID 是 token 在模型内部的整数表示。除非需要特别区分，本章在描述 LLM 运算
  时会交替使用 token 和 token ID 这两个术语。

  \item \textbf{Embedding（嵌入）层}把分词器生成的 token（即 token ID）转换成嵌入
  向量，也就是 token 的浮点向量表示。这里交替使用 embedding 和 embedding vector
  这两个说法。转换通常通过查找嵌入表完成。训练过程中也会从文档中学习每个 token 的
  嵌入值。之所以要进行这种转换，是因为 LLM 是使用数值方法训练的机器学习模型。训练
  通常采用梯度下降方法，要求模型函数可微，并且输入和参数采用连续形式；浮点表示是
  对连续值的一种近似。

  \item \textbf{位置编码（positional encoding）层}把序列顺序信息加入输入 token 的
  嵌入向量。位置编码告诉注意力机制 token 在序列中的物理位置和相互距离；这些因素会
  在判断此前输入和输出 token 对下一个 token 生成的相关性时发挥作用。

  当前对话中的完整 token 序列（包括用户输入、系统添加的输入和模型输出）会按照生成
  顺序累积成一个矩阵。矩阵的每一行是一个 token 的嵌入向量，该矩阵作为第一个
  Transformer 层的输入。

  \item \textbf{Transformer 层}对嵌入进行变换，以提取 token 之间的依赖关系。每个
  Transformer 层包含一个多头注意力子层、一个加法与归一化子层、一个前馈子层，以及
  另一个加法与归一化子层。

  \item \textbf{多头注意力子层}由并行排列的多个注意力头组成。每个注意力头通过矩阵
  乘法对嵌入执行线性投影。线性投影所用矩阵是在训练过程中学习得到的，它们实际上
  对嵌入向量执行坐标变换，使后续的相似度比较更有效。投影结果再经过另一次矩阵乘法，
  测量一个 token 与此前已经生成的所有 token 之间的相似度。多个注意力头的输出会被
  拼接起来，通常再经过线性变换，形成多头注意力子层的组合输出。

  \item \textbf{加法与归一化（Add \& Norm）子层}位于多头注意力子层之后。每个前馈
  子层或注意力子层之后的加法操作都建立残差连接，以保留原始输入的信息。归一化操作
  通过把结果向量归一化到下一层所需的数值范围，稳定计算过程并改善收敛。归一化通常
  使用输入嵌入向量元素的均值 $\mu$、标准差 $\sigma$ 以及学习得到的参数 $\gamma$、
  $\beta$。

  \item \textbf{前馈子层}通常由两个线性变换和中间的 ReLU 激活组成。直观地说，注意力
  层识别 token 之间的相似性或依赖关系，并根据相似度加权组合此前 token 的嵌入向量，
  生成新的嵌入版本；前馈子层则为这些新的嵌入版本增加语义含义。关于 ReLU 和激活
  函数的定义，参见附录 B。

  \item 另一个 \textbf{加法与归一化子层}位于前馈子层之后，用来保留前馈子层输入中的
  信息、改善收敛，并最终生成该 Transformer 层的输出。
\end{enumerate}

模型中 Transformer 层的数量称为模型深度，不同 LLM 的深度差别很大。执行大规模矩阵
乘法的注意力头，是 Transformer 架构中计算开销最大的部分，也是本章后续内容的重点。
最后，\textbf{反嵌入（un-embedding）层}把最后一个 Transformer 层输出的嵌入转换成
token 上的概率分布，用于生成下一个输出 token。

模型维度对于生成高质量回答非常重要。第一维是模型深度，即 Transformer 层的数量。
更深的模型提供更多层级的分类器，因而能够捕获更一般的依赖模式。第二维是每个嵌入
向量中的元素数量，称为头维度（head dimension）。更宽的嵌入向量能够捕获更多语义
信息，区分不同概念，并发现更有意义的依赖关系。第三维是每个注意力子层中的头数量；
不同注意力头能够捕获 token 之间不同类型的依赖。模型参数量约等于这些模型维度相关
规模的乘积，实际模型可能包含数十亿甚至数万亿个参数（模型权重）。

虽然 LLM 需要较长上下文才能生成准确响应，但用户输入通常过于简短且隐含信息较多，
不能提供充分的上下文。一个常见补救方法是使用检索增强生成（retrieval-augmented
generation，RAG）：检索相关文档，把更具体、更及时的上下文信息与用户输入拼接起来，
形成更大、更丰富的上下文。

在 LLM 训练期间，最重要的计算通常是大规模矩阵--矩阵乘法。注意力头和前馈子层都会
执行这种运算，并可以使用第 15 章介绍的高级优化策略。正因为矩阵乘法占据主要计算
开销，注意力头成为 LLM 中最昂贵的计算组件。下一节将详细讨论注意力头的设计及其所
执行的矩阵乘法。

\section{多头注意力}
\label{sec:ch20-multi-head-attention}

每个注意力子层包含多个头。多个头共同执行多个分类任务，并同时捕获不同类型的依赖，
例如位置依赖（token 之间的距离）、语义依赖（同义词或主题）、上下文依赖（把冠词与
名词联系起来）和句法依赖（把动词与宾语联系起来）。这种注意力头的专门化是在训练
过程中形成的。不同头的权重矩阵采用不同初始化；反向传播的梯度下降会避免不同头
之间的功能重叠，并通过 dropout（随机置零输出或概率）和层归一化（归一化激活值）等
技术抑制过拟合。

注意力头彼此独立，因此可以很容易地分配到多个线程块上并发计算。不同头的设计相同，
所以下面只展示一个注意力头。图 \ref{fig:20-2} 给出了一个注意力头的计算过程。
输入矩阵 $X$ 的每一行对应 LLM 当前已经生成的 token 的嵌入向量，包括系统提示词、
用户输入、系统对用户输入的补充信息以及已经生成的输出 token。对于第一个 Transformer
层的注意力子层，$X$ 的每一行包含位置编码组件生成的嵌入；对于后续 Transformer 层，
输入 $X$ 就是前一层的输出。

当前已经生成的 token 数量，也就是矩阵 $X$ 的行数，称为序列长度或上下文长度，后文
统一记作 $N$。在聊天机器人应用中，随着用户与系统的对话持续进行，$N$ 会逐渐增长。
支持较大 $N$ 的模型可以利用更多上下文来消除词语和短语的歧义、判断上下文相关性，
并保持连续回答之间的一致性，因此通常能够生成更可靠的 token。注意力子层通常遵循
因果策略：当前 token 的生成只能依赖此前生成的 token，而忽略未来 token。

每个嵌入向量中的元素数量，即头维度，记作 $d$。因此，$X$ 是一个 $N\times d$ 矩阵。
头维度和每个子层中的头数量通常被称为 LLM 的隐藏维度，它们是模型的内部设计参数，
通常不会直接暴露给 LLM 用户。

\begin{figure}[htbp]
  \centering
  \small
  \begin{tabular}{c@{\quad}c@{\quad}c@{\quad}c}
    $X_{N\times d}$ & $\xrightarrow{\;W_Q\;}$ & $Q_{N\times d}$ &
    $\xrightarrow{\quad}$\\[0.35em]
    \rotatebox{90}{$\begin{matrix}\xrightarrow{\;W_K\;}\\[-0.2em]
    \xrightarrow{\;W_V\;}\end{matrix}$} & &
    $\begin{matrix}K_{N\times d}\\ V_{N\times d}\end{matrix}$ &\\[0.5em]
    \multicolumn{4}{c}{$S=QK^{\mathsf T}\ (N\times N)
      \quad\xrightarrow{\;+\ M\;}\quad
      P=\operatorname{softmax}(S+M)\ (N\times N)
      \quad\xrightarrow{\;\times V\;}\quad
      O_{N\times d}$}
  \end{tabular}
  \caption{一个注意力头计算中涉及的矩阵（依据原书图 20.2 临时重绘）。
  为简化图示，省略了作用于 $S=QK^{\mathsf T}$ 的缩放因子 $1/\sqrt d$。}
  \label{fig:20-2}
\end{figure}

注意力子层的第一步是计算三个矩阵：query（查询，$Q$）、key（键，$K$）和 value
（值，$V$）。如图 \ref{fig:20-2} 左侧所示，这些矩阵由 $X$ 的行向量经过线性投影
或坐标变换得到，即把输入矩阵分别乘以三个权重矩阵 $W_Q$、$W_K$ 和 $W_V$。在本章
的简化讨论中，这三个坐标变换矩阵都是 $d\times d$ 矩阵，并由训练过程得到。$Q$、
$K$ 和 $V$ 的每一行中的元素，都是 $X$ 同一行中 $d$ 个元素的线性组合；三种变换的
差异由 $W_Q$、$W_K$ 和 $W_V$ 中的权重值决定。

\noindent\textbf{译者注：}为简化说明，本章假设三个权重矩阵是方阵。一般情况下，
$W_Q$、$W_K$ 和 $W_V$ 的尺寸分别可以是 $d_x\times d_q$、$d_x\times d_k$ 和
$d_x\times d_v$；由于需要计算 $QK^{\mathsf T}$，必须有 $d_q=d_k$。

注意力子层的第二步由下面的公式给出，也对应图 \ref{fig:20-2} 的右侧：
\begin{equation}
  O =
  \operatorname{softmax}\!\left(
    \frac{QK^{\mathsf T}}{\sqrt{d}} + M
  \right)V .
\label{eq:ch20-attention}
\end{equation}

由于 $Q$ 是 $N\times d$ 矩阵，$K^{\mathsf T}$ 是 $d\times N$ 矩阵，矩阵乘法
$QK^{\mathsf T}$ 会生成一个 $N\times N$ 矩阵，记作图 \ref{fig:20-2} 中的 $S$。
$QK^{\mathsf T}$ 的元素 $(i,j)$ 是 token $i$ 的线性投影嵌入向量（$Q$ 的第 $i$ 行）
与 token $j$ 的线性投影嵌入向量（$K^{\mathsf T}$ 的第 $j$ 列，也就是 $K$ 的第 $j$
行）之间的内积。如果两个输入向量都已归一化，内积可以看作它们的余弦相似度。两个
向量越相似，它们之间的夹角越小，余弦相似度也越大。因此，$N\times N$ 的乘积矩阵
可以看作嵌入向量之间的相似度地图。

缩放因子 $1/\sqrt d$ 作用于 $S=QK^{\mathsf T}$ 的所有元素，用于防止头维度较大时
矩阵乘法的点积过度增长。过大的 $S$ 元素会使 softmax 饱和，并在基于梯度下降的训练
中产生很小、甚至消失的梯度 \cite{ch20-1}。图 \ref{fig:20-2} 及相关讨论为简洁起见
省略了缩放因子；稍后可以看到，把它并入 CUDA softmax 实现非常直接。

掩码矩阵 $M$ 会把 $QK^{\mathsf T}/\sqrt d$ 中所有列索引大于行索引的元素加上
$-\infty$。这样便实现了因果策略：生成某个 token 时，不能受到该 token 之后才生成
的 token 的影响。

softmax 函数把 $QK^{\mathsf T}+M$ 的每一行转换为概率分布。设该矩阵第 $r$ 行、第
$c$ 列的元素为 logit $l_{r,c}$，softmax 生成的概率为 $p_{r,c}$。为了避免指数函数
溢出，令 $m_r$ 为该行的最大元素，并采用数值稳定的形式：
\begin{equation}
  p_{r,c}
  =
  \frac{e^{l_{r,c}}}{\sum_{j=1}^{N} e^{l_{r,j}}}
  =
  \frac{e^{l_{r,c}-m_r}}
       {\sum_{j=1}^{N}e^{l_{r,j}-m_r}},
  \qquad
  m_r=\max_{1\leq j\leq N} l_{r,j}.
\label{eq:ch20-softmax}
\end{equation}

直接计算 $e^{l_{r,c}}$ 可能溢出，因此在取指数之前先从 $l_{r,c}$ 中减去
$m_r=\max_{1\leq j\leq N}(l_{r,j})$。这个操作保持了数值稳定性，同时不改变结果：
从 $l_{r,c}$ 中减去 $m_r$ 等价于把分子和分母都乘以 $e^{-m_r}$。

由于因果掩码，行索引 $r$ 小于列索引 $c$ 的 logit 会变成 $-\infty$，于是
$e^{l_{r,c}-m_r}=0$。因此，矩阵
$P=\operatorname{softmax}(QK^{\mathsf T}+M)$ 的每一行，都是 token $r$ 的嵌入与
此前生成 token 的嵌入之间相似度所决定的概率分布。softmax 输出矩阵是一个下三角
矩阵：当 $r<c$ 时，元素 $p_{r,c}$ 为零。

注意力子层的最后一步是计算 $N\times N$ 的概率矩阵 $P$ 与 $V$ 的矩阵乘法。对于
每个 token，这个矩阵乘法都会生成一个新的嵌入向量；新向量的每个元素，都是 $V$
矩阵中同一列元素的相似度/概率加权和。换句话说，输出矩阵 $O$ 中每个 token 的新
嵌入，是此前生成 token 的嵌入向量的相似度加权和。之后，$O$ 会经过加法与归一化、
前馈等子层，成为下一个注意力子层的输入矩阵 $X$。

\begin{figure}[htbp]
  \centering
  \small
  \begin{tabular}{c@{\qquad}c}
    \fbox{\begin{minipage}{0.39\textwidth}
      \centering
      \textbf{摘要/预填充阶段}\\[0.3em]
      系统提示词 + 用户输入 + RAG 文档\\
      $\downarrow$ tokenizer / embedding\\
      $X$：一整个 token 序列\\
      $\downarrow$\\
      一次 Transformer 前向传播\\
      $\downarrow$\\
      第一个输出 token
    \end{minipage}}
    &
    \fbox{\begin{minipage}{0.39\textwidth}
      \centering
      \textbf{生成/解码阶段}\\[0.3em]
      把新 token 追加到上下文\\
      $\downarrow$\\
      每次处理一个新 token\\
      $\downarrow$\\
      生成下一个 token\\
      $\circlearrowleft$ 直到 EOS
    \end{minipage}}
  \end{tabular}
  \caption{LLM 推理的摘要（prefill）阶段和生成（decode）阶段
  （依据原书图 20.3 临时重绘）。}
  \label{fig:20-3}
\end{figure}

基于图 \ref{fig:20-1} 的仅解码器架构进行 LLM 推理时，首先进入摘要阶段：系统提示词、
用户输入以及系统补充的 token（例如 RAG 检索结果）拼接起来，形成第一个 Transformer
层的输入矩阵 $X$。摘要阶段中，每个注意力头涉及的矩阵和计算如图 \ref{fig:20-2}
所示。累积的 $X$ 依次经过各个 Transformer 层；每层生成下一层的输入。最后一层完成
后，反嵌入层生成一个文本片段形式的输出 token。下一次推理迭代进入自回归阶段，也称
解码阶段或生成阶段。新生成的文本片段被送回分词器，其嵌入向量作为 $X$ 的新行追加
到第一个 Transformer 层，然后再次激活各层以生成下一个文本片段。

图 \ref{fig:20-3} 用一个玩具例子说明两个阶段。摘要阶段首先把系统提示和用户提示
组成的初始序列送入分词器和嵌入层，并累积到矩阵 $X$ 中。为简化起见，例子省略系统
提示词和用户输入中通常存在的前导语。用户输入为 “The title of this book is
Programming”，希望 LLM 补全本书标题。

假设这个例子使用完整单词作为 token。摘要阶段中，分词器、嵌入层和位置编码层把用户
输入产生的七个 token 的嵌入向量累积到第一个 Transformer 层的 $X$ 矩阵中。LLM
随后调用第一个 Transformer 层；该层完成后激活第二层，以此类推。每层被激活时都会
计算自己的 $K$、$V$、$Q$、$QK^{\mathsf T}$ 和 $O$，然后激活下一层。所有 Transformer
层和最后的反嵌入层完成后，LLM 生成第一个输出 token “Massively”。

生成第一个输出 token 后，LLM 进入生成阶段。每个输出 token 都根据初始提示以及此前
所有迭代生成的输出 token 产生。在例子中，摘要阶段末尾生成的 “Massively” 会成为
生成阶段第一次迭代的输入，并由这次迭代产生下一个输出 token。对于每个新输出 token，
LLM 都会生成一个对应的新嵌入向量，也就是矩阵 $X$ 的新行，并让它经过各个 Transformer
层。第一次生成迭代得到 “Parallel”，它会成为下一次迭代的输入。生成阶段会持续进行，
直到产生 end-of-sequence（EOS）token。

\section{在 CUDA 中实现注意力}
\label{sec:ch20-cuda-attention}

图 \ref{fig:20-1} 所示 Transformer 架构中的注意力计算包含两次矩阵乘法、一次矩阵
逐元素缩放、一次矩阵加法和 softmax 计算。矩阵乘法可以使用自定义 CUDA 矩阵乘法
内核，也可以调用 cuBLAS 等库中高度优化的 GEMM 函数。缩放因子 $1/\sqrt d$ 可以
在自定义内核执行 $QK^{\mathsf T}$ 时应用，也可以作为 cuBLAS GEMM API 的缩放参数
$\alpha$。注意力计算中唯一的新部分是 softmax 的实现，因此本节重点讨论 CUDA
softmax，并把其余部分留作练习。

图 \ref{fig:20-4} 展示了一个把与 $M$ 的矩阵加法和 softmax 计算融合在同一内核中的
CUDA 实现。融合内核避免在 softmax 之前额外遍历输入矩阵 $S$，从而提高算术强度和
GPU 利用率。

\begin{figure}[htbp]
  \centering
  \begin{lstlisting}[language=C++,basicstyle=\ttfamily\tiny]
// Figure 20.4: causal softmax kernel
#define BLOCK_SIZE 256
#define NEG_INFINITY (-INFINITY)

__global__ void softmax_kernel(float* S, float* D, float* P, int N)
{
    typedef cub::BlockReduce<float, BLOCK_SIZE> BlockReduce;
    __shared__ typename BlockReduce::TempStorage temp_store;
    __shared__ float max_or_sum;

    float* S_row = &S[blockIdx.x * N];
    float max_val_thread = NEG_INFINITY;

    for (int idx = threadIdx.x;
         idx <= blockIdx.x;
         idx += blockDim.x) {
        if (S_row[idx] > max_val_thread)
            max_val_thread = S_row[idx];
    }

    __syncthreads();
    float max_val_row =
        BlockReduce(temp_store).Reduce(max_val_thread, cub::Max());

    if (threadIdx.x == 0)
        max_or_sum = max_val_row;
    __syncthreads();
    max_val_row = max_or_sum;

    float sum_thread = 0.f;
    for (int idx = threadIdx.x;
         idx <= blockIdx.x;
         idx += blockDim.x) {
        sum_thread += exp(S_row[idx] - max_val_row);
    }

    __syncthreads();
    float sum_row =
        BlockReduce(temp_store).Reduce(sum_thread, cub::Sum());

    if (threadIdx.x == 0)
        max_or_sum = sum_row;
    __syncthreads();
    sum_row = max_or_sum;

    for (int idx = threadIdx.x;
         idx < N;
         idx += blockDim.x) {
        P[blockIdx.x * N + idx] =
            idx <= blockIdx.x
                ? exp(S_row[idx] - max_val_row) / sum_row
                : 0.f;
    }

    if (threadIdx.x == 0)
        D[blockIdx.x] = sum_row;
}
  \end{lstlisting}
  \caption{融合因果掩码和 softmax 的 CUDA 内核（依据原书图 20.4 整理）。
  代码标识符保持原文；本图为可复制版本。}
  \label{fig:20-4}
\end{figure}

该内核的执行配置是一个一维线程块网格。内核假定线程块数量
\code{gridDim.x} 等于矩阵 $S$ 的行数，也就是序列长度 $N$。底本把未使用的
\code{gridDim.y} 和 \code{gridDim.z} 描述为 0；在合法的 CUDA 启动配置中，未使用
维度通常取 1，本译稿把这句话理解为“只使用 x 维”，而不把 0 当作可直接传入的网格
维度。这个假设把上下文长度限制在 x 维允许的最大线程块数以内。若需要更长上下文，
可以在内核中增加外层循环，让网格中的线程块迭代处理 $S$ 的所有行。

每个线程块是一维块，包含 \code{BLOCK_SIZE} 个线程。由于一个线程块协同处理 $S$
的一行，应选择合适的 \code{BLOCK_SIZE}：既要减少线程块处理该行所需的迭代次数，
也要尽量提高 GPU 上 SM 的占用率。

当缩放因子已经在 $QK^{\mathsf T}$ 矩阵乘法中应用，并且因果策略已经融合到 softmax
中时，一行 $S$ 的每个元素都可以看作公式 \eqref{eq:ch20-softmax} 中的
$l_{r,c}$。代码前半部分先找出分配给线程块的那一行的最大元素，即公式中的 $m_r$。
每个线程检查自己负责的一段元素，并把局部最大值保存在私有变量
\code{max_val_thread} 中。循环条件 \code{idx <= blockIdx.x} 很关键：它忽略行索引
\code{blockIdx.x} 小于列索引 \code{idx} 的元素，从而不显式构造掩码矩阵 $M$ 就实现
因果策略。

线程退出循环后，每个线程只持有自己检查部分的最大值。块内同步确保所有线程都完成
检查，然后调用 CUB 的 \code{BlockReduce} 求整行最大值。归约结果由
\code{threadIdx.x == 0} 的线程获得；该线程把结果写入共享内存变量
\code{max_or_sum}，再通过块内同步把结果广播给所有线程。代码的这部分本质上是带
线程粗化的归约。

接下来，代码以相同的方式求公式 \eqref{eq:ch20-softmax} 的分母。线程首先计算私有
部分和；循环条件继续实现因果策略。同步后，CUB 的 \code{BlockReduce} 求出整行的
指数和，线程 0 再通过共享内存广播最终结果。最后一个循环遍历行中的全部元素，把
因果掩码区域写成 0，把其余元素写入输出矩阵 $P$。分母同时写入输出向量 $D$，以便
训练阶段复用。

总之，该 softmax 内核在一个内核调用中同时实现因果策略和 softmax，只使用块级
\code{__syncthreads()}，不需要额外的全局屏障。读者可以在此基础上补充生成 $S$
和 $O$ 的矩阵乘法，完成整个注意力子层。

\section{KV 缓存}
\label{sec:ch20-kv-cache}

如果按照图 \ref{fig:20-3} 的直接方式实现，图 \ref{fig:20-1} 的 Transformer 架构
存在明显低效之处。在生成阶段的每次迭代中，每个 Transformer 层都要执行五次矩阵
乘法。由于序列长度 $N$ 可能增长到数万个 token 甚至更多，这些矩阵乘法的开销会非常
大。但从一次解码迭代到下一次迭代，输入矩阵 $X$ 只增加一行，因此有必要问：每次
迭代真的都需要执行完整的矩阵乘法吗？答案是否定的。

\begin{figure}[htbp]
  \centering
  \small
  \resizebox{0.98\textwidth}{!}{%
  \begin{tabular}{c@{\qquad}c@{\qquad}c}
    \begin{tabular}{c}
      $X_N$\\[-0.2em]
      \chTwentyCell{x_0}\\
      \chTwentyCell{x_1}\\
      \vdots\\
      \chTwentyNewCell{x_N}
    \end{tabular}
    &
    $\xrightarrow{\;W_Q,W_K,W_V\;}$
    &
    \begin{tabular}{c}
      $Q_N,\ K_N,\ V_N$\\[-0.2em]
      前 $N$ 行保持不变\\
      新增最后一行
    \end{tabular}\\[1em]
    \multicolumn{3}{c}{
      \begin{tabular}{c}
        $QK^{\mathsf T}$ 从 $N\times N$ 变为 $(N+1)\times(N+1)$\\[-0.1em]
        \chTwentyTextCell{旧} \chTwentyTextCell{旧} \chTwentyTextCell{新}\\
        \chTwentyTextCell{旧} \chTwentyTextCell{旧} \chTwentyTextCell{新}\\
        \chTwentyTextCell{新} \chTwentyTextCell{新} \chTwentyTextCell{新}
      \end{tabular}
    }
  \end{tabular}
  }%
  \caption{从一次 LLM 迭代到下一次迭代时减少矩阵乘法工作的机会
  （依据原书图 20.5 临时重绘；灰色表示新增或需要处理的部分）。}
  \label{fig:20-5}
\end{figure}

图 \ref{fig:20-5} 突出了从迭代 $i$（$N=i$）到迭代 $i+1$（$N=i+1$）时各矩阵
新增的部分。输入矩阵 $X$ 在底部增加一个新行，即新 token 的嵌入向量，因此尺寸从
$N\times d$ 变成 $(N+1)\times d$。$Q$、$K$ 和 $V$ 的情况类似：它们都变成
$(N+1)\times d$ 矩阵，而新增的底部行分别由 $X$ 的新底行与 $W_Q$、$W_K$、$W_V$
执行向量--矩阵乘法得到。

生成 $QK^{\mathsf T}$ 稍微复杂一些。新矩阵的尺寸是 $(N+1)\times(N+1)$，但除第
$(N+1)$ 行和第 $(N+1)$ 列外，其余内容基本与上一次迭代相同。例如，乘积矩阵的
元素 $(0,0)$ 由 $Q$ 的第 0 行与 $K$ 的第 0 行（即 $K^{\mathsf T}$ 的第 0 列）
的内积生成，而这两行在两次迭代中都没有变化。因此该元素也不变。一般而言，只要
$r$ 和 $c$ 都小于 $N$，$QK^{\mathsf T}$ 的元素 $(r,c)$ 就由两次迭代中均未变化
的 $Q$、$K$ 行生成，因而保持不变。

新 $N$ 行的 $QK^{\mathsf T}$ 可以通过新 $Q$ 行与完整 $K^{\mathsf T}$ 的向量--矩阵
乘法计算。新 $N$ 列中的元素由于因果策略的零掩码都为零，因为此前生成的 token 不
应该受到后来生成 token 的影响；唯一例外是新行新列交叉处的元素，它属于新行计算
的一部分。

softmax 的输出也可以从上一次迭代增量更新。softmax 作用于 $QK^{\mathsf T}$ 的
每一行；除了新行之外，每一行新增的第 $N$ 个元素都是 0，因此不会稀释原有元素的
概率，也不会改变这些行的 softmax 结果。读者可以验证：在迭代 $N+1$ 中，$O$ 的
前 $N$ 行仍然与上一次迭代相同，只需用新一行 softmax 输出与新 $V$ 做向量--矩阵
乘法即可得到新行。

\begin{figure}[htbp]
  \centering
  \small
  \begin{tabular}{c@{\quad}c@{\quad}c}
    \begin{tabular}{c}
      新输入行\\
      $X'$\\[0.3em]
      $\downarrow$\\
      $Q'$、$K'$、$V'$
    \end{tabular}
    &
    $\longrightarrow$
    &
    \begin{tabular}{c}
      KV cache 中的旧 $K,V$\\
      $\left[\begin{array}{c}K\\V\end{array}\right]
      \quad+\quad
      \left[\begin{array}{c}K'\\V'\end{array}\right]$\\[0.4em]
      $\begin{gathered}
        \downarrow\\[-0.1em]
        S'=Q'[K;K']^{\mathsf T}\\
        P'=\operatorname{softmax}(S'+M)\\
        O'=P'[V;V']
      \end{gathered}$
    \end{tabular}
  \end{tabular}
  \caption{使用 KV cache 的生成阶段涉及的向量和矩阵
  （依据原书图 20.6 临时重绘）。}
  \label{fig:20-6}
\end{figure}

图 \ref{fig:20-6} 展示了生成 $Q$、$K$、$V$ 和 $O$ 新部分时涉及的向量和矩阵。每次
迭代，多头注意力子层只需要接收 $X$ 的一个新行，因此只需向下一层传递 $O$ 的一个
新行。不过，生成阶段的注意力计算仍然需要完整的 $K$ 和 $V$。一种常见优化是把
$K$ 和 $V$ 记忆化，也就是保存起来供以后复用，这便形成了 KV cache。
\footnote{这里的 KV 与键值存储（key-value store）文献中的术语不应混淆，但二者
具有类比关系：LLM 推理输出的是 $V$ 的加权和，权重与 query 和 key 的接近程度
相关。}

\begin{figure}[htbp]
  \centering
  \small
  \begin{tabular}{c}
    \fbox{\begin{minipage}{0.82\textwidth}
      \centering
      \textbf{Prefill（摘要）}\\
      初始 token 序列 $\longrightarrow$ 各 Transformer 层\\
      $K,V$ 矩阵 $\longrightarrow$ 每层自己的 KV cache
    \end{minipage}}\\[0.7em]
    $\Downarrow$\\[-0.2em]
    \fbox{\begin{minipage}{0.82\textwidth}
      \centering
      \textbf{Decode/Generation（生成）}\\
      新 token 的嵌入 $\longrightarrow$ 新 $Q',K',V'$\\
      $K',V'$ 追加到 cache $\longrightarrow$ 计算 $O'$ $\longrightarrow$ 下一个 token
    \end{minipage}}
  \end{tabular}
  \caption{使用 KV cache 的 LLM 摘要与生成阶段
  （依据原书图 20.7 临时重绘）。}
  \label{fig:20-7}
\end{figure}

图 \ref{fig:20-7} 在图 \ref{fig:20-3} 的基础上，展示了带 KV cache 的两阶段推理。
为了避免每次迭代重复计算，每个 Transformer 层生成的 $K$、$V$ 都会写入该层的
KV cache，供以后生成 token 时复用。

在摘要阶段，所有 Transformer 层生成初始的 $K$、$V$ 矩阵并填充自己的 KV cache。
因此，带 KV cache 的摘要阶段也称为预填充（prefill）阶段：各层在这一阶段把初始
$K$、$V$ 写入 cache。随后进入生成阶段，每个输出 token 都基于最后一个输出 token
以及此前所有迭代中的 $K$、$V$ 矩阵生成。最近生成的文本片段先产生一个新的嵌入
向量，并送入第一个 Transformer 层。每层先执行三个向量--矩阵乘法，分别生成 $Q$、
$K$、$V$ 的一行，即图 \ref{fig:20-6} 中的 $Q'$、$K'$、$V'$。LLM 把新行追加到
KV cache 中原有的 $K$、$V$，形成当前迭代的完整 $K$、$V$，再计算 $O'$，并把它
送入该层的后续子层，生成下一层的新输入向量。

预填充阶段仍然需要对大量 token 执行一次 Transformer 传递，并调用多次大规模矩阵--
矩阵乘法（GEMM）。它通常受计算能力限制，算术强度较高，能够较充分地利用 GPU。
下一节介绍的 FlashAttention 可以通过减少矩阵乘法与 softmax 之间中间结果的全局
内存流量，进一步加速预填充阶段。

相比之下，生成阶段每生成一个输出 token 就要执行一次 Transformer 传递。使用 KV
cache 后，每次传递主要执行向量--矩阵乘法（GEMV），通常受内存带宽限制，算术强度
较低，因而容易使 GPU 的计算资源利用不足。第 20.6 节将介绍 batching 和 speculative
decoding 等提高生成阶段算术强度的优化方法。

\section{FlashAttention}
\label{sec:ch20-flash-attention}

如果像图 \ref{fig:20-4} 那样使用独立的 softmax 内核，那么每个内核边界都会产生
全局屏障，并且整个矩阵需要多次从全局内存加载或写回。全局屏障和全局内存访问会
成为重要的性能瓶颈；FlashAttention 公式化方法 \cite{ch20-6} 可以缓解这一问题。

FlashAttention 对注意力头进行数学上的重新表达，使每个注意力头中的运算可以重新
排序并融合到一个内核中，同时以分块方式组织。在融合内核中，如图 \ref{fig:20-8}
所示，每个线程块取出 $Q$ 的一个水平面板（连续的若干行）、完整的 $K^{\mathsf T}$
和完整的 $V$，计算它们对 $O$ 的一个水平面板的贡献。图中用阴影表示属于某个线程
块的输入和输出。

\begin{figure}[htbp]
  \centering
  \small
  \resizebox{0.98\textwidth}{!}{%
  \begin{tabular}{c@{\qquad}c@{\qquad}c}
    \begin{tabular}{c}
      $Q\ (N\times d)$\\
      \fbox{\begin{tabular}{cccc}
        \chTwentyNewCell{q}&\chTwentyNewCell{q}&\chTwentyNewCell{q}&\chTwentyNewCell{q}\\
        \chTwentyNewCell{q}&\chTwentyNewCell{q}&\chTwentyNewCell{q}&\chTwentyNewCell{q}\\
        \chTwentyCell{q}&\chTwentyCell{q}&\chTwentyCell{q}&\chTwentyCell{q}
      \end{tabular}}\\[-0.2em]
      $Q_i:\ B_r\times d$
    \end{tabular}
    &
    \begin{tabular}{c}
      $K^{\mathsf T}\ (d\times N)$\\
      \fbox{\begin{tabular}{cccc}
        \chTwentyNewCell{k}&\chTwentyCell{k}&\chTwentyCell{k}&\chTwentyCell{k}\\
        \chTwentyNewCell{k}&\chTwentyCell{k}&\chTwentyCell{k}&\chTwentyCell{k}\\
        \chTwentyNewCell{k}&\chTwentyCell{k}&\chTwentyCell{k}&\chTwentyCell{k}
      \end{tabular}}\\[-0.2em]
      $K^{\mathsf T}_j:\ d\times B_c$
    \end{tabular}
    &
    \begin{tabular}{c}
      $V\ (N\times d)$\\
      \fbox{\begin{tabular}{cccc}
        \chTwentyNewCell{v}&\chTwentyNewCell{v}&\chTwentyNewCell{v}&\chTwentyNewCell{v}\\
        \chTwentyCell{v}&\chTwentyCell{v}&\chTwentyCell{v}&\chTwentyCell{v}\\
        \chTwentyCell{v}&\chTwentyCell{v}&\chTwentyCell{v}&\chTwentyCell{v}
      \end{tabular}}\\[-0.2em]
      $V_j:\ B_c\times d$
    \end{tabular}\\[1em]
    \multicolumn{3}{c}{
      $\displaystyle
      Q_iK^{\mathsf T}_j\longrightarrow S_i
      \xrightarrow{\ \operatorname{softmax}\ }P_i
      \longrightarrow P_iV_j\longrightarrow O_i$}\\[0.8em]
    \multicolumn{3}{c}{
      \fbox{\begin{tabular}{c}
        $Q_i,\ O_i$：寄存器分块\qquad
        $K^{\mathsf T}_j,\ S_i/P_i,\ V_j$：共享内存分块\\
        每次只处理 $B_c$ 列/行，并在线合并 softmax 状态
      \end{tabular}}}
  \end{tabular}
  }%
  \caption{FlashAttention 的分块方式。共享内存中的同一个 $S_i$ tile 先保存
  $S$，再原地复用为 $P$（依据原书图 20.8 临时重绘）。}
  \label{fig:20-8}
\end{figure}

此外，FlashAttention 允许线程块以分块方式遍历 $K^{\mathsf T}$ 和 $V$，同时计算
它们对 $O$ tile 的贡献。这些 tile 足够小，可以放入片上内存。每次迭代，线程块加载
$K^{\mathsf T}$ 的一个垂直面板（连续列）和 $V$ 的一个水平面板（连续行），计算它们
对 $O$ tile 的贡献。所有线程块可以并发执行，不需要全局屏障。这样，中间矩阵只以
tile 的形式保留在片上内存中，不必把它们写入全局内存，从而消除了中间矩阵对应的
全局内存流量。

FlashAttention 是对 Transformer 注意力机制的精确重构，与原始公式在数学上完全等价，
而不是近似算法。它既适用于训练（前向和反向传播），也适用于推理（前向传播）。本节
只讨论前向传播。反向传播在概念上更简单，因为不需要 softmax 重缩放，但它需要五次
矩阵乘法，而前向传播只需要两次；关于反向传播的细节，可参考原始 FlashAttention
论文 \cite{ch20-6}。

FlashAttention 的关键是分块。然而，注意力头中的分块不像矩阵乘法中的分块那样直接，
因为公式 \eqref{eq:ch20-softmax} 要求对输入的每一行执行多次扫描。softmax 计算依赖
于 $QK^{\mathsf T}+M$ 第 $r$ 行的最大值 $m_r$，以及生成概率时作为分母使用的
\[
  D_r=\sum_{j=1}^{N}e^{l_{r,j}-m_r}.
\]
按照这个形式，至少需要先扫描一行以同时收集最大值和分母，再进行另一次扫描来计算
输出矩阵 $P$ 的元素。为了保证整个 $QK^{\mathsf T}+M$ 矩阵已经准备好，softmax
内核之前还需要全局同步。

在 CUDA 中，可以通过结束一个内核并启动另一个内核实现全局屏障；也可以启动持久化
网格，使用 cooperative groups 的网格级同步 API，从而避免内核调用开销。不过，适合
矩阵乘法的线程网格配置通常不同于适合 softmax 的配置，这也是基线 CUDA 实现采用
不同内核、分别配置矩阵乘法和 softmax 的重要原因。无论选择哪种全局同步策略，多次
扫描 $S$ 的整行都会引入大量全局内存流量，显著限制吞吐量。

幸运的是，可以为分母定义部分级数，消除 softmax 对整行执行多次扫描的需要。令
$A$ 为第 $r$ 行中列位置的一个子集，定义：
\begin{equation}
  D_{r,A}=\sum_{j\in A}e^{l_{r,j}-m_{r,A}},
\label{eq:ch20-partial-denominator}
\end{equation}
其中 $m_{r,A}$ 是子集 $A$ 中元素的最大值。当 $A$ 覆盖该行的所有列位置时，
$D_{r,A}$ 就等于 $D_r$。

如果能够推导出一个组合规则，把 $A$ 的部分级数与另一个列子集 $B$ 的部分级数
合并成 $A\cup B$ 的部分级数，就可以独立地（甚至并行地）计算各个子集，然后在
结果准备好时逐步合并它们 \cite{ch20-7}。当合并后的子集覆盖第 $r$ 行的所有列时，
就能得到正确的 $D_r$。把 $D_{r,A}$ 与 $D_{r,B}$ 合并的规则为：
\begin{equation}
\begin{aligned}
  D_{r,A\cup B}
  &=\sum_{j\in A\cup B}e^{l_{r,j}-m_{r,A\cup B}}\\
  &=D_{r,A}e^{m_{r,A}-m_{r,A\cup B}}
    +D_{r,B}e^{m_{r,B}-m_{r,A\cup B}},
\end{aligned}
\label{eq:ch20-denominator-composition}
\end{equation}
其中 $m_{r,A\cup B}$ 是 $m_{r,A}$ 和 $m_{r,B}$ 中较大的一个。上式的拆分要求
$A$ 和 $B$ 不包含重复元素。在下面的 CUDA 实现中，$A$ 和 $B$ 是 $S$ 每一行的
不同 tile，因此这个条件天然满足。只要同时保存每个子集对应的最大值，就可以使用
这个规则合并两个部分和。

FlashAttention 进一步把同样的组合规则应用到输出矩阵 $O$。先考虑 $O$ 的一个元素
$o_{r,c}$，它是 $P$ 的第 $r$ 行与 $V$ 的第 $c$ 列之间的内积。只考虑行 $r$ 的
子集 $A$ 时，定义部分输出：
\begin{equation}
  o_{r,c,A}
  =\sum_{j\in A}p_{r,j,A}v_{j,c}
  =\frac{1}{D_{r,A}}
    \sum_{j\in A}e^{l_{r,j}-m_{r,A}}v_{j,c}.
\label{eq:ch20-partial-output}
\end{equation}
$p_{r,j,A}$ 是只考虑子集 $A$ 时的中间概率值，$m_{r,A}$ 也只取子集 $A$ 中
元素的最大值，而不是整行的最大值。随着更多元素并入 $A$，当 $A$ 最终覆盖整行时，
$p_{r,j,A}$ 就会变成完整 softmax 中的 $p_{r,j}$。

根据分母的组合规则，部分输出的组合规则为：
\begin{equation}
\begin{aligned}
  o_{r,c,A\cup B}
  &=
  \frac{D_{r,A}e^{m_{r,A}-m_{r,A\cup B}}}
       {D_{r,A\cup B}}\,o_{r,c,A}\\
  &\quad+
  \frac{D_{r,B}e^{m_{r,B}-m_{r,A\cup B}}}
       {D_{r,A\cup B}}\,o_{r,c,B}.
\end{aligned}
\label{eq:ch20-output-composition}
\end{equation}

也就是说，在把两个子集的部分输出合并为完整输出前，必须根据各自的最大值和分母
对它们进行重缩放。这个规则对于行的任意子集都成立，但当子集是由 $Q$ 和
$K^{\mathsf T}$ 的 tile 生成的连续片段时尤其有用。图 \ref{fig:20-8} 正是利用
这一规则实现 FlashAttention 的分块方法 \cite{ch20-8}。

\begin{figure}[htbp]
  \centering
  \begin{lstlisting}[language=C++,basicstyle=\ttfamily\tiny]
#define WARP_SIZE 32
#define BLOCK_SIZE 512
#define N_WARPS (BLOCK_SIZE / WARP_SIZE)
#define B_r 32
#define B_r_warp (B_r / N_WARPS)
#define B_c 32
#define d 128

__global__ void flashattention_forward_kernel(
    const float* Q,
    const float* K,
    const float* V,
    const int N,
    const float scaling,
    float* out_D,
    float* out_O)
{
    const int T_r = N / B_r;
    const int T_c = N / B_c;
    const int d_size = d > WARP_SIZE ? d / WARP_SIZE : 1;

    __shared__
    float KT_j[B_c * d + ((B_c * d) >> LOG_NUM_BANKS)];
    __shared__ float S_i[B_r][B_c];
    __shared__ float V_j[B_c][d];

    for (int i = blockIdx.x; i < T_r; i += gridDim.x) {
        float O_i[B_r_warp][d_size];
        float D_i[B_r_warp];
        float m_i[B_r_warp];

        initialize(O_i, D_i, m_i);

        float Q_i[B_r_warp][d_size];
        load_Q(Q, Q_i, i);

        for (int j = 0; j < T_c; j++) {
            load_KT_and_V(K, V, KT_j, V_j, j);
            __syncthreads();

            for (int ii = 0; ii < B_r_warp; ii++) {
                float curr_max_warp =
                    compute_S_and_max(
                        Q_i, KT_j, S_i, i, j, ii, scaling);

                float last_m =
                    update_m_and_D(
                        m_i, D_i, curr_max_warp, ii);

                float curr_sum_warp =
                    compute_P_and_update_D(
                        S_i, m_i, D_i, i, j, ii);

                compute_O(
                    S_i, V_j, O_i, ii, last_m, m_i);
            }
        }

        __syncthreads();
        store_O(out_O, out_D, O_i, D_i, i, N);
    }
}
  \end{lstlisting}
  \caption{FlashAttention 前向传播的 CUDA 内核主体（依据原书图 20.9 整理）。
  该示例处理一个注意力头；多头版本可以增加网格维度。}
  \label{fig:20-9}
\end{figure}

图 \ref{fig:20-9} 的内核计算一个注意力头的前向传播。完整网格参与该计算；在更
完整的实现中，可以通过增加线程块网格的其他维度，在同一个内核中计算多个头。相同
的分块方式也适用于 LLM 推理的预填充阶段。图中示例使用 32 位浮点数，而当前高性能
实现通常使用 16 位甚至 8 位浮点数，并利用 GPU Tensor Core 的吞吐能力。

该内核假定输入矩阵 $Q$、$K$、$V$ 都是 $N\times d$，主机代码负责分配全局内存、
初始化输入，并为输出矩阵 $O$（代码中的 \code{out_O}）分配空间。内核还生成一个
长度为 $N$ 的输出向量 $D$（代码中的 \code{out_D}），其中每个元素是 $P$ 对应行
的最终 softmax 分母。这个向量主要供训练阶段使用；推理阶段不一定需要它，但保留
它可以使同一个内核同时用于训练。

在高层次上，FlashAttention 内核使用一维线程块网格，每个块负责生成 $O$ 的一个水平
面板。输出 tile 的水平尺寸是注意力头宽度 $d$，垂直尺寸由 $B_r$ 指定。$B_r$ 的
选择受片上内存容量和充分利用 GPU 所需线程块数的共同限制，目标是最大化片上数据
复用。每个 warp 负责的行数为 $B_{r,\mathrm{warp}}$，即 $B_r$ 除以线程块中的
warp 数量。

内核最外层循环保证所有 $O$ 面板都被生成。每次迭代中，一个线程块处理一个 $Q$
面板并生成对应的 $O$ 面板；所有线程块共同生成 \code{gridDim.x} 个面板，直到
生成全部 $T_r=N/B_r$ 个面板。

对于一个 $Q$ 面板，线程块分三步工作：首先把 $B_r\times d$ 的 $Q$ 水平面板与
完整的 $K^{\mathsf T}$ 相乘，得到 $B_r\times N$ 的 $S$ 水平面板；其次对 $S$ 应用
softmax，得到相应的 $P$ 面板；最后把 $P$ 与完整的 $V$ 相乘，得到 $B_r\times d$
的 $O$ tile。朴素实现可能用三个内核完成这三步，需要把 $S$ 和 $P$ 写入并重新从
全局内存读取，导致低算术强度和额外的内核调用开销。公式 \eqref{eq:ch20-output-composition}
允许这三步在一个内核中以迭代方式融合，并在片上内存中逐块完成，从而不必在全局
内存中实体化 $S$ 和 $P$。

这里同时使用共享内存分块和寄存器分块：$Q$、$O$ tile 放在寄存器中，$K^{\mathsf T}$、
$S/P$ 和 $V$ tile 放在共享内存中。代码中的 \code{S_i} 先保存 $S$，随后原地复用
为 softmax 输出 $P$。每个线程还在寄存器中保存行最大值 $m_i$、softmax 分母
$D_i$ 和输出 $O_i$；这些状态由 \code{initialize()} 初始化。$Q_i$ 则由
\code{load_Q()} 从全局内存加载。

一个线程只持有这些 tile 的一部分元素，并通过 warp shuffle 原语与同一个 warp 中
的线程交换数据。$Q$ 的面板进一步划分为大小为 $B_{r,\mathrm{warp}}\times d$ 的
子面板，每个子面板分配给一个 warp。下面逐个解释各个 device function。

\noindent\textbf{译者注：}这里的 $Q$ 分配方式对应 FlashAttention-2
\cite{ch20-8}。早期版本把 $K$ 和 $V$ 分配给不同 warp（称为 split-K），需要把
中间结果写入共享内存、同步并归约，效率较低；本节示例让 warp 之间不交换中间数据，
只在同一 warp 内交换元素。

\begin{figure}[htbp]
  \centering
  \begin{lstlisting}[language=C++,basicstyle=\ttfamily\scriptsize]
__device__ inline void load_Q(
    const float* Q,
    float Q_i[B_r_warp][d_size],
    int i)
{
    for (int ii = 0; ii < B_r_warp; ii++) {
        for (int dd = laneIdx(), ddd = 0;
             dd < d;
             dd += WARP_SIZE, ddd++) {
            Q_i[ii][ddd] =
                Q[(B_r * i +
                   B_r_warp * warpIdx() +
                   ii) * d + dd];
        }
    }
}
  \end{lstlisting}
  \caption{FlashAttention：把 $Q$ 加载到寄存器中的 device function
  （依据原书图 20.10 整理）。}
  \label{fig:20-10}
\end{figure}

在每个子面板中，行 $ii$ 以交错方式分配给 warp 中的线程。也就是说，一行被划分为
长度为 \code{WARP_SIZE} 的区段，每个区段的元素按顺序分配给 warp 中的线程。这样，
warp 可以以合并方式加载每个区段。每个线程最终在寄存器中保存
$B_{r,\mathrm{warp}}\times(d/\mathrm{WARP\_SIZE})$ 个 $Q$ 元素。示例为简洁起见
假设 $d$ 是 \code{WARP_SIZE} 的倍数；若要支持任意 $d$，需要额外处理边界。

暂时忽略 softmax，可以更容易地理解融合内核中的两个矩阵乘法。线程块先把自己的
$Q$ tile 加载到寄存器，并调用 \code{initialize()} 把 $O$ tile 的元素初始化为零，
然后以分块方式迭代执行两个矩阵乘法。这个迭代过程有两个层次：块级和 warp 级。
块级的每次迭代内部又是一个 warp 级迭代。

在块级迭代中，线程块从 $K^{\mathsf T}$ 的左端和 $V$ 的顶部开始；对
$K^{\mathsf T}$ 从左向右、对 $V$ 从上向下遍历。每次处理 $K^{\mathsf T}$ 的
$B_c$ 列和 $V$ 的 $B_c$ 行。由于 $K^{\mathsf T}$ 有 $N$ 列、$V$ 有 $N$ 行，
线程块需要 $T_c=N/B_c$ 次迭代。每次迭代都要合并 $O$ tile 的所有元素，因此
$B_c$ 应尽量大，以减少合并次数；但它又受可用共享内存限制。示例要求 $B_c$ 是
\code{WARP_SIZE} 的倍数，从而不需要边界检查或填充；更通用的实现可以放宽这一假设。

\begin{figure}[htbp]
  \centering
  \begin{lstlisting}[language=C++,basicstyle=\ttfamily\scriptsize]
__device__ inline void load_KT_and_V(
    const float* K,
    const float* V,
    float* KT_j,
    float V_j[B_c][d],
    int j)
{
    for (int jj = 0; jj < B_c; jj++) {
        for (int dd = threadIdx.x;
             dd < d;
             dd += blockDim.x) {
            KT_j[addr(dd * B_c + jj)] =
                K[(B_c * j + jj) * d + dd];

            V_j[jj][dd] =
                V[(B_c * j + jj) * d + dd];
        }
    }
}
  \end{lstlisting}
  \caption{FlashAttention：把 $K^{\mathsf T}$ 和 $V$ 加载到共享内存中的 device
  function（依据原书图 20.11 整理）。}
  \label{fig:20-11}
\end{figure}

在块级迭代开始时，线程块把 $K^{\mathsf T}$ 的 $B_c$ 列和 $V$ 的 $B_c$ 行加载到
共享内存。为了提高效率，$K^{\mathsf T}$ 从 $K$ 中以转置方式加载。为避免共享内存
bank 冲突，图 \ref{fig:20-9} 的声明为 $K^{\mathsf T}_j$ 的每一列增加了 padding。
因此访问 $K^{\mathsf T}_j$ 时，要把线性下标 $x$ 转换为第 16 章介绍的
\[
  \operatorname{addr}(x)=x+\left(x\mathbin{\mathord{\gg}}\mathrm{LOG\_NUM\_BANKS}\right).
\]
图 \ref{fig:20-9} 中的块级屏障保证所有 tile 都已经加载完成，才会开始第一次矩阵
乘法。

每个 $O$ tile 的生成由嵌套循环实现。每个 warp 取其 $Q$ 子面板的一行，与当前
$K^{\mathsf T}$ 面板执行向量--矩阵乘法，把结果写入 $S_i$ 的对应行；随后使用
softmax 把 $S_i$ 原地转换成 $P$ 的一段；最后把该行与当前 $V$ 面板执行向量--
矩阵乘法，把贡献合并到自己的 $O$ 子面板。

\begin{figure}[htbp]
  \centering
  \begin{lstlisting}[language=C++,basicstyle=\ttfamily\tiny]
__device__ inline float compute_S_and_max(
    const float Q_i[B_r_warp][d_size],
    const float* KT_j,
    float S_i[B_r][B_c],
    int i,
    int j,
    int ii,
    float scaling)
{
    typedef cub::WarpReduce<float> WarpReduce;
    __shared__ typename WarpReduce::TempStorage temp_store;

    float curr_max = NEG_INFINITY;

    for (int jj = laneIdx(); jj < B_c; jj += WARP_SIZE) {
        float S_ij = 0.f;

        for (int dd = 0; dd < d; dd++) {
            float q =
                __shfl_sync(0xFFFFFFFF,
                            Q_i[ii][dd / WARP_SIZE],
                            dd % WARP_SIZE);
            S_ij += q * KT_j[addr(dd * B_c + jj)];
        }

        int row = B_r * i + B_r_warp * warpIdx() + ii;
        int col = B_c * j + jj;

        S_ij = row < col ? NEG_INFINITY : scaling * S_ij;
        S_i[B_r_warp * warpIdx() + ii][jj] = S_ij;

        if (S_ij > curr_max)
            curr_max = S_ij;
    }

    float curr_max_warp =
        cub::WarpReduce<float>(temp_store).Reduce(
            curr_max, cub::Max());

    curr_max_warp =
        __shfl_sync(0xFFFFFFFF, curr_max_warp, 0);

    return curr_max_warp;
}
  \end{lstlisting}
  \caption{FlashAttention：计算 $S$ tile 及其最大值的 device function
  （依据原书图 20.12 整理）。}
  \label{fig:20-12}
\end{figure}

函数 \code{compute_S_and_max()} 把当前 $K^{\mathsf T}$ 面板进一步划分为每段
\code{WARP_SIZE} 列。循环中的每个线程计算当前 $Q$ 行与当前 $K^{\mathsf T}$ 子面板
一列的内积。由于 $Q$ 的行元素以交错方式分布在 warp 线程的寄存器中，线程需要
通过 \code{__shfl_sync()} 把当前所需的 $Q$ 元素广播给同一 warp 的其他线程。
这样，一个 warp 每次处理一个 $d\times\mathrm{WARP\_SIZE}$ 的 $K^{\mathsf T}$
子面板，并产生 $S_i$ 对应行中的一段。

代码随后比较全局矩阵 $S$ 中的行、列索引：如果列索引大于行索引，就把对应共享
内存位置写成 $-\infty$；否则把点积乘以缩放因子 $1/\sqrt d$ 后写入 $S_i$。最后，
每个线程求自己计算元素的局部最大值，再使用 CUB 的 \code{WarpReduce} 求 warp
内的最大值，并用 shuffle 把归约结果广播给 warp 中的所有线程。

\noindent\textbf{补充说明：}这里的示例只展示了共享内存和寄存器分块；实际实现还
可以采用层次化分块、Tensor Core 等更复杂的优化方法。

\begin{figure}[htbp]
  \centering
  \begin{lstlisting}[language=C++,basicstyle=\ttfamily\scriptsize]
__device__ inline void compute_O(
    const float S_i[B_r][B_c],
    const float V_j[B_c][d],
    float O_i[B_r_warp][d_size],
    int ii,
    float last_m,
    float* m_i)
{
    for (int dd = laneIdx(), ddd = 0;
         dd < d;
         dd += WARP_SIZE, ddd++) {
        O_i[ii][ddd] *= exp(last_m - m_i[ii]);

        float O_ij = 0.f;
        for (int jj = 0; jj < B_c; jj++) {
            O_ij +=
                S_i[B_r_warp * warpIdx() + ii][jj]
                * V_j[jj][dd];
        }

        O_i[ii][ddd] += O_ij;
    }
}
  \end{lstlisting}
  \caption{FlashAttention：计算 $O$ tile 贡献的 device function
  （依据原书图 20.13 整理）。}
  \label{fig:20-13}
\end{figure}

\code{compute_O()} 中，每个 warp 把自己的 $S_i$ 行与 $V$ 的当前子面板相乘，
生成对 $O$ tile 的部分贡献。一个 $O$ 元素实际上是完整 $S$ 行与完整 $V$ 列的
内积；当前 tile 只贡献其中连续的 $B_c$ 项，因此每次块级迭代都要把结果累加到
寄存器中的 $O_i$。

加入 softmax 后，每次生成 $S_i$ 行，warp 中的线程先共同求出该行的最大值 $m_i$
和分母 $D_i$。随后线程块根据这两个状态得到 $P$ 行，把 $S_i$ 原地转换成 $P$，
再把 $P$ 与 $V$ 当前子面板相乘，并按照公式 \eqref{eq:ch20-output-composition}
合并到 $O$ tile。

\begin{figure}[htbp]
  \centering
  \begin{lstlisting}[language=C++,basicstyle=\ttfamily\scriptsize]
__device__ inline float update_m_and_D(
    float* m_i,
    float* D_i,
    float curr_max_warp,
    int ii)
{
    float last_m = m_i[ii];
    float D = D_i[ii];

    if (m_i[ii] < curr_max_warp) {
        m_i[ii] = curr_max_warp;
        D *= exp(last_m - curr_max_warp);
    }

    D_i[ii] = D;
    return last_m;
}
  \end{lstlisting}
  \caption{FlashAttention：更新 $m$ 和 $D$ 的 device function
  （依据原书图 20.14 整理）。}
  \label{fig:20-14}
\end{figure}

图 \ref{fig:20-14} 中，线程先把当前行最大值 $m_i$ 保存到 \code{last\_m}，如果
当前 tile 的最大值 \code{curr\_max\_warp} 更大，就更新累计最大值。这里的三个值
分别对应公式 \eqref{eq:ch20-denominator-composition} 和
\eqref{eq:ch20-output-composition} 中的 $m_{r,A}$、$m_{r,B}$ 和
$m_{r,A\cup B}$：$A$ 是此前已经处理的列子集，$B$ 是当前新 tile。更新最大值
后，代码对已有分母进行相应的指数重缩放。

\begin{figure}[htbp]
  \centering
  \begin{lstlisting}[language=C++,basicstyle=\ttfamily\tiny]
__device__ inline float compute_P_and_update_D(
    float S_i[B_r][B_c],
    float* m_i,
    float* D_i,
    int i,
    int j,
    int ii)
{
    typedef cub::WarpReduce<float> WarpReduce;
    __shared__ typename WarpReduce::TempStorage temp_store;

    float curr_sum = 0.f;

    for (int jj = laneIdx(); jj < B_c; jj += WARP_SIZE) {
        int row = B_r * i + B_r_warp * warpIdx() + ii;
        int col = B_c * j + jj;

        float P_ij =
            row < col
                ? 0.f
                : exp(S_i[B_r_warp * warpIdx() + ii][jj]
                      - m_i[ii]);

        S_i[B_r_warp * warpIdx() + ii][jj] = P_ij;
        curr_sum += P_ij;
    }

    float curr_sum_warp =
        cub::WarpReduce<float>(temp_store).Reduce(
            curr_sum, cub::Sum());

    curr_sum_warp =
        __shfl_sync(0xFFFFFFFF, curr_sum_warp, 0);

    D_i[ii] += curr_sum_warp;
    return curr_sum_warp;
}
  \end{lstlisting}
  \caption{FlashAttention：计算 $P$ 并更新 $D$ 的 device function
  （依据原书图 20.15 整理）。}
  \label{fig:20-15}
\end{figure}

\code{compute_P_and_update_D()} 完成 softmax 的指数步骤，并再次应用因果策略。它
把结果写回共享内存中的 \code{S_i}，从而原地把 $S$ tile 转换成 $P$ tile，同时
累积当前 tile 对分母的贡献。CUB 的 \code{WarpReduce} 求出当前 warp 的部分和，
shuffle 把结果广播给所有线程，然后更新 \code{D_i[ii]}。

每个 $o_{r,c}$ 都应该接收 $P$ 行中所有元素的贡献，但当前 $P$ tile 与 $V$ tile
的矩阵乘法只包含该内积的一段。因此，在把当前矩阵乘法结果累加到输出 tile 之前，
必须依据公式 \eqref{eq:ch20-output-composition} 缩放已有的 $O$ 内容和当前的
矩阵乘法结果。图 \ref{fig:20-13} 中的第一行先把已有 $O_i$ 乘以
$\exp(\texttt{last\_m}-m_i)$，对应组合公式的第一项；随后内层循环把当前
$P$/$V$ tile 的贡献加入。当前新贡献已经使用截至当前 tile 的最大值计算，因此
不需要再次重缩放。

\begin{figure}[htbp]
  \centering
  \begin{lstlisting}[language=C++,basicstyle=\ttfamily\scriptsize]
__device__ inline void store_O(
    float* out_O,
    float* out_D,
    const float O_i[B_r_warp][d_size],
    const float D_i[B_r_warp],
    int i,
    int N)
{
    for (int ii = 0; ii < B_r_warp; ii++) {
        int row =
            B_r * i + B_r_warp * warpIdx() + ii;

        for (int dd = 0; dd < d_size; dd++) {
            int col = dd * WARP_SIZE + laneIdx();

            if (row < N && col < d)
                out_O[row * d + col] =
                    O_i[ii][dd] / D_i[ii];
        }

        if (laneIdx() == 0)
            out_D[row] = D_i[ii];
    }
}
  \end{lstlisting}
  \caption{FlashAttention：把 $O$ 和分母写回全局内存的 device function
  （依据原书图 20.16 整理）。}
  \label{fig:20-16}
\end{figure}

线程块完成迭代后，输出 tile 已经包含概念上的完整 $S$、$P$ 矩阵贡献，尽管这些
矩阵从未完整生成，只以 tile 的形式存在于共享内存中。写回时，\code{store_O()}
把寄存器中的 $O_i$ 和 softmax 分母 $D_i$ 写入全局内存。写入 $O$ 前必须除以最终
分母，这正是公式 \eqref{eq:ch20-output-composition} 中的归一化步骤。

\noindent\textbf{译者注：}上述代码是用于解释算法的数据布局示例，不应直接视为
生产级 FlashAttention 实现。它假设 $N$ 能被 $B_r$ 和 $B_c$ 整除，$B_c$ 与 $d$
满足 warp 对齐，并依赖第 16 章中定义的 \code{addr()}、\code{laneIdx()} 和
\code{warpIdx()}。实际工程还需要处理任意尺寸、不同精度、边界条件、共享内存
临时存储生命周期，以及 CUB 归约对象的安全复用问题。

本节描述的 FlashAttention 内核适用于不带 KV cache 的基线 Transformer，也适用于
带 KV cache 推理的预填充阶段。在带缓存的生成（decode）阶段，单个头中的 $Q$、$S$、
$P$ 和 $O$ 矩阵会变成向量 $Q'$、$S'$、$P'$ 和 $O'$；矩阵--矩阵乘法也会变为
向量--矩阵乘法。为了保持与图 \ref{fig:20-8} 类似的工作分配方式，可以把多个
请求，即多个 $Q_i$，batch 成一个矩阵。

上面介绍的迭代过程只是众多可能实现中的一种。如果使用 Tensor Core 执行每个 warp
中的两次矩阵乘法，就可能需要采用不同的数据分配和迭代组织方式；读者可以继续探索
这些实现变体。

较新的 FlashAttention-3 \cite{ch20-9} 还利用 Hopper GPU 架构的新能力进一步提升
性能，包括：Tensor Memory Accelerator（TMA），它可以异步地把 tile 从全局内存搬到
共享内存供 Tensor Core 使用而不占用计算线程；Tensor Core 的 WarpGroup-wide
Matrix Multiply Accumulate（WGMMA）指令；以及 Tensor Core 的 FP8 精度。通过 warp
专门化，这一版本能够更好地重叠数据移动和计算，并用异步的分块 GEMM 隐藏 softmax
执行。

\section{KV cache 的算术强度与内存需求}
\label{sec:ch20-kv-cache-ai}

图 \ref{fig:20-2} 说明了 LLM 推理预填充阶段中 $K$ 和 $V$ 矩阵的生成。$K$、$V$
的元素会被计算出来并缓存到内存中；在生成阶段（图 \ref{fig:20-6}），每生成一个
新 token，就会生成 $K$ 和 $V$ 的一行，并把它追加到 KV cache。这些新行通过向量--
矩阵运算生成，算术强度较低，容易使 GPU 的计算资源利用不足。因此，batching 和
speculative decoding 等优化方法都试图提高生成阶段的算术强度。

Batching 是所有深度神经网络中都很常见的优化方法，传统上用于提高计算强度。例如，
一个有 1024 个输入、4096 个输出的线性层，如果使用 16 位精度且 batch size 为 1，
算术强度只有每字节 1 个浮点运算（FLOPS/B）；当 batch size 增加到 512 时，算术
强度可以提高到 315 FLOPS/B。算术强度的提高来自不同输入向量之间对权重的复用，
而高算术强度对于充分利用 GPU 十分重要。

不过，LLM 在 batch size 和上下文长度（也就是 KV cache 大小）之间存在权衡。多个
用户的输入查询通常使用同一个模型，因此可以共享模型权重，提高 GPU 在不同计算
阶段中的利用率。图 \ref{fig:20-17} 左侧示意了在生成 $Q$、$K$、$V$ 的全连接层
中进行 batching 的方式；batching 可以使 QKV 投影的算术强度近似提高 batch size
倍数。然而，KV cache 依赖每个用户自己的提示词和上下文，注意力层仍然必须维护并
访问各自独立的 cache。因此，LLM 中的 batching 会受到模型权重和 KV cache 内存
需求的共同限制。

\begin{figure}[htbp]
  \centering
  \small
  \begin{tabular}{c@{\qquad}c}
    \fbox{\begin{minipage}{0.40\textwidth}
      \centering
      \textbf{共享模型权重}\\[0.3em]
      请求 1 的 $X_1$ \\
      请求 2 的 $X_2$ \\
      请求 3 的 $X_3$\\[0.3em]
      $\left.\begin{array}{c}X_1\\X_2\\X_3\end{array}\right\}
      \xrightarrow{\quad W_Q,W_K,W_V\quad}
      \left\{\begin{array}{c}Q_1,K_1,V_1\\
      Q_2,K_2,V_2\\Q_3,K_3,V_3\end{array}\right.$
    \end{minipage}}
    &
    \fbox{\begin{minipage}{0.40\textwidth}
      \centering
      \textbf{独立 KV cache}\\[0.3em]
      请求 1 $\longrightarrow$ $\mathrm{KVCache}_1$\\
      请求 2 $\longrightarrow$ $\mathrm{KVCache}_2$\\
      请求 3 $\longrightarrow$ $\mathrm{KVCache}_3$\\[0.3em]
      注意力计算必须访问各自上下文
    \end{minipage}}
  \end{tabular}
  \caption{LLM 推理中的 batching：QKV 投影可以共享权重，而注意力仍需为不同
  会话维护独立 KV cache（依据原书图 20.17 临时重绘）。}
  \label{fig:20-17}
\end{figure}

下面分析带 KV cache 和 batching 的 LLM 推理内存需求。主要贡献者有两个：模型权重
（例如 $W_Q$、$W_K$ 和 $W_V$）以及 KV cache。模型权重的大小由模型参数数量决定，
通常达到数十亿量级。例如，一个包含 70 亿参数、以 16 位格式（FP16 或 BF16）加载的
模型大约需要 14 GB。关于模型权重内存需求的更多细节，可参考 \cite{ch20-10,ch20-11}。

KV cache 的大小取决于 Transformer 层数 $l$ 和隐藏维度。隐藏维度是 query 头数量
$h_q$ 与嵌入维度 $d$ 的乘积。对于多头注意力（multi-head attention，MHA），每个
token 的 KV cache 大小（字节数）为：
\begin{equation}
  \mathrm{KVCacheSize}_{\mathrm{MHA}}
  =2\ell h_q d p .
\label{eq:ch20-kv-mha-per-token}
\end{equation}
这里乘以 2 是因为同时保存 $K$ 和 $V$，而 $p$ 是每个元素所占的字节数。

实际场景还要计算 batch 中所有序列的所有 token 的 cache 大小。对于 batch 中有 $b$
个、长度均为 $N$ 的序列，总 KV cache 大小为：
\begin{equation}
  \mathrm{TotalKVCacheSize}_{\mathrm{MHA}}
  =bN\cdot 2\ell h_q d p .
\label{eq:ch20-kv-mha-total}
\end{equation}

例如，在 16 位精度下，考虑长度为 4096 的单个序列。对于 Google 的 PaLM 2，
$\ell=64$、隐藏维度 $h_qd=8192$，总 KV cache 大小为
$1\times4096\times2\times64\times8192\times2=8$ GB。对于 OpenAI 的 GPT-3，
相应大小为
$1\times4096\times2\times96\times12288\times2=18$ GB。对于 Meta 的 Llama 2 7B，
大小为
$1\times4096\times2\times32\times4096\times2=2$ GB。如此大的 KV cache 很容易使
批量推理的内存需求超过单个 GPU 的容量，因此 LLM 推理通常需要多 GPU 甚至多节点
系统。

上述分析假设 batch 中所有序列长度相同，但现实中不同用户的输入序列和生成输出
序列通常长度不同。这样，batching 可能造成线程块之间的工作量不均衡，导致 GPU
资源利用率下降。为避免这种情况，先进的 LLM 服务系统会使用 in-flight batching
机制 \cite{ch20-12}：系统一次调度并执行多个请求，并在某个请求完成后，立即把新请求
安排到当前工作量最低的 batch 中。

虽然 batching 能有效提高线性投影操作的算术强度，但对推理阶段的 MHA 并不十分有效。
可以通过近似数据移动量和浮点运算量来估算它的算术强度 \cite{ch20-13}。考虑一个
层、长度为 $N$ 的序列，并令 $p=1$。根据公式 \eqref{eq:ch20-kv-mha-per-token}，
KV cache 大小为 $2h_qdN$；读入 $Q'$ 和写出 $O'$ 的数据移动量为 $2h_qd$；计算
$(Q'K^{\mathsf T})V$ 的浮点运算量为 $2Nh_qd$，这里忽略 softmax。于是，MHA 的
算术强度近似为：
\begin{equation}
  \mathrm{AI}_{\mathrm{MHA}}
  \approx
  \frac{2Nh_qd}{2h_qd+2h_qdN}
  =\frac{N}{1+N}
  \approx 1 .
\label{eq:ch20-ai-mha}
\end{equation}
因此，MHA 在生成阶段的算术强度可能低至每字节约 1 个浮点运算。

\noindent\textbf{译者注：}这里的 batching 增益是理论估计；实际算术强度还取决于
实现所采用的 tile 尺寸和数据布局。

提高 LLM 推理算术强度的另一种方法是 speculative decoding（推测解码）
\cite{ch20-14}。这是一种模型服务优化：一个很小的 draft model（草稿模型）先预测
大型 target model（目标模型）的若干输出 token；目标模型随后验证这些预测。虽然
目标模型单次延迟更高，但它可以并行验证草稿模型一次生成的一批 token。

\begin{figure}[htbp]
  \centering
  \small
  \begin{tabular}{c}
    \fbox{\begin{minipage}{0.82\textwidth}
      \centering
      \textbf{Draft model}\\
      已接受的上下文 $\longrightarrow$
      \texttt{Programming}\quad\texttt{Massively}\quad\texttt{Efficient}
    \end{minipage}}\\[0.6em]
    $\Downarrow\ \text{并行验证}$\\[-0.1em]
    \fbox{\begin{minipage}{0.82\textwidth}
      \centering
      \textbf{Target model}\\
      \texttt{Programming}\;（接受）
      \qquad \texttt{Massively}\;（接受）
      \qquad \texttt{Efficient}\;（拒绝）\\[0.3em]
      接受前两个 token，拒绝第三个，然后开始下一轮草稿预测
    \end{minipage}}
  \end{tabular}
  \caption{推测解码：草稿模型预测多个 token，目标模型并行验证
  （依据原书图 20.18 临时重绘）。}
  \label{fig:20-18}
\end{figure}

图 \ref{fig:20-18} 给出了一个推测解码例子。推测深度为 3，即草稿模型连续预测三个
输出 token，目标模型再并行验证它们。在第一轮中，目标模型接受 “Programming”
和 “Massively”，但拒绝 “Efficient”。目标模型验证结束后，被接受的 token 会并入
下一轮草稿模型的输入，草稿模型再预测三个新 token。

推测解码与 batching 相似，因此也能提高线性投影操作的算术强度。不过，它是一种特殊
形式的 batching：批中的所有 token 共享一组初始 token。因此，在注意力计算中，这些
共同初始 token 对应的 KV cache 数据只需加载一次，就可以复用于所有输出 token。这种
KV cache 数据复用正是算术强度提高的原因；没有推测解码时，同一份 cache 数据需要
为每个输出 token 分别加载。

\section{降低注意力机制的内存需求}
\label{sec:ch20-reducing-attention-memory}

如前所述，注意力机制使用多个头，每个头使用不同的权重矩阵来生成 $Q$、$K$ 和 $V$
的投影。这是多头注意力（MHA）内存需求较大的主要原因之一。为了在保持可接受模型
准确率的同时降低内存需求，人们提出了多查询注意力（multi-query attention，MQA）
和分组查询注意力（grouped-query attention，GQA）两种方法。图 \ref{fig:20-19}
以简化方式展示 MHA、MQA、GQA 和多头潜在注意力（multi-head latent attention，
MLA）的内存组织。

\begin{figure}[htbp]
  \centering
  \small
  \begin{tabular}{c@{\qquad}c}
    \fbox{\begin{minipage}{0.37\textwidth}
      \centering
      \textbf{MHA}\\[0.25em]
      $Q_1\ Q_2\ Q_3\ Q_4$\\
      $K_1\ K_2\ K_3\ K_4$\\
      $V_1\ V_2\ V_3\ V_4$\\
      每个头都有独立的 $K,V$
    \end{minipage}}
    &
    \fbox{\begin{minipage}{0.37\textwidth}
      \centering
      \textbf{MQA}\\[0.25em]
      $Q_1\ Q_2\ Q_3\ Q_4$\\
      $\quad K_{\!*}$\\
      $\quad V_{\!*}$\\
      所有头共享一组 $K,V$
    \end{minipage}}\\[1em]
    \fbox{\begin{minipage}{0.37\textwidth}
      \centering
      \textbf{GQA}\\[0.25em]
      $Q_1\ Q_2\ |\ Q_3\ Q_4$\\
      $K_1\quad\ |\ K_2$\\
      $V_1\quad\ |\ V_2$\\
      每组 query 头共享 $K,V$
    \end{minipage}}
    &
    \fbox{\begin{minipage}{0.37\textwidth}
      \centering
      \textbf{MLA}\\[0.25em]
      $Q_1\ Q_2\ Q_3\ Q_4$\\
      $\downarrow$ 低秩联合压缩\\
      $\quad z_{\mathrm{latent}}$\\
      需要时再上投影恢复 $K,V$
    \end{minipage}}
  \end{tabular}
  \caption{MHA、MQA、GQA 和 MLA 中 $Q$、$K$、$V$ 的共享关系
  （依据原书图 20.19 临时重绘）。}
  \label{fig:20-19}
\end{figure}

MQA 在多个注意力头之间共享 key 和 value，从而减小隐藏维度。MQA 的总 KV cache
大小为：
\begin{equation}
  \mathrm{TotalKVCacheSize}_{\mathrm{MQA}}
  =bN\cdot 2\ell d p .
\label{eq:ch20-kv-mqa}
\end{equation}
与公式 \eqref{eq:ch20-kv-mha-total} 相比，缓存大小按头数量的倒数缩减。由于 KV
cache 减小，MQA 的算术强度近似提高到 MHA 的 $h_q$ 倍：
\begin{equation}
  \mathrm{AI}_{\mathrm{MQA}}
  \approx
  \frac{2Nh_qd}{2h_qd+2dN}
  =\frac{Nh_q}{h_q+N}
  \approx h_q .
\label{eq:ch20-ai-mqa}
\end{equation}

GQA \cite{ch20-15} 在 MHA 和 MQA 之间进行折中：把 query heads 划分成少量的
$g_q$ 个组，同组 query 头共享 key 和 value。GQA 的总内存需求为：
\begin{equation}
  \mathrm{TotalKVCacheSize}_{\mathrm{GQA}}
  =bN\cdot 2\ell\cdot\frac{h_q}{g_q}d p .
\label{eq:ch20-kv-gqa}
\end{equation}
GQA 的算术强度近似提高到组数 $g_q$：
\begin{equation}
  \mathrm{AI}_{\mathrm{GQA}}
  \approx
  \frac{2Nh_qd}
       {2h_qd+2\frac{h_q}{g_q}dN}
  =\frac{Nh_q}{h_q+\frac{h_q}{g_q}N}
  \approx g_q .
\label{eq:ch20-ai-gqa}
\end{equation}

MHA、GQA 和 MQA 都可以使用第 20.5 节介绍的 FlashAttention 分块方法。三种方案
在计算量、内存需求和模型准确率之间有不同的权衡。

尽管 MQA 和 GQA 能够降低 KV cache 的内存需求，但有限的 GPU 内存仍然会限制 batch
size 和上下文长度。更糟糕的是，KV cache 通常会按照最大可能输入预留空间，导致
内存碎片和浪费。PagedAttention \cite{ch20-12} 是一种尝试解决这一问题的方法，
其思想来自操作系统的分页机制：把 KV cache 划分成大小相同的块，并在需要计算时，
从更大的系统内存把所需块加载到 GPU 内存中。

管理 KV cache 的另一种方法是多头潜在注意力（MLA）\cite{ch20-16}。MLA 对 $K$ 和
$V$ 执行低秩联合压缩，大幅降低内存使用。其核心思想是：把每个 token 在所有头中
对应的 $K$、$V$ 向量压缩成一个潜在向量，如图 \ref{fig:20-19} 所示。压缩在所有
头和所有层中进行，因此 KV cache 的内存需求可以显著降低：
\begin{equation}
  \mathrm{TotalKVCacheSize}_{\mathrm{MLA}}
  =bN\ell d p .
\label{eq:ch20-kv-mla}
\end{equation}

由于多个 query 头共享一个潜在头，MLA 的算术强度会提高。底本给出的公式为：
\begin{equation}
  \mathrm{AI}_{\mathrm{MLA}}
  \approx
  \frac{2Nh_qd}{2h_qd+dN}
  =\frac{2Nh_q}{h_q+N}
  \approx 2h_q .
\label{eq:ch20-ai-mla}
\end{equation}

\noindent\textbf{译者注：}按照首行分母严格约分，中间项应为
$2Nh_q/(2h_q+N)$；底本排版给出的中间式是 $2Nh_q/(h_q+N)$。两种形式在
$N$ 很大时都趋近于 $2h_q$，本译稿保留底本等式并标注这一代数化简疑点，未将其
悄然改写。

压缩方案所使用的投影矩阵在训练过程中学习得到。推理时，模型会在使用 $K$ 和 $V$
进行注意力计算之前，把潜在表示乘以两个上投影矩阵，以恢复（解压）$K$ 和 $V$。
因此，迭代之间只需要在各注意力头的 KV cache 中保存压缩后的潜在表示；只有当某个
头在生成阶段处于活动状态时，才解压出该头所需的 $K$ 和 $V$。

扩展 LLM 的另一种正交技术是专家混合（mixture of experts，MoE），例如 Mixtral
\cite{ch20-17}。MoE 是一种稀疏激活架构：门控机制会根据每个输入决定使用少量
专门化子网络中的哪一部分。这样，模型可以拥有极大的参数量，同时降低每个 token
实际需要执行的计算量。

\section{小结}
\label{sec:ch20-summary}

本章介绍了支撑大语言模型（LLM）的基础运算。由于主流 LLM 大多建立在 Transformer
架构之上，本章重点讨论 Transformer 特有的计算组件，即注意力头。我们首先从朴素的
多内核实现开始，重点分析注意力头中最具新意的 softmax 内核；随后讨论 KV cache，
这是减少 LLM 生成阶段每次迭代中冗余计算的重要优化。接着介绍 FlashAttention 的
公式化方法及其注意力头内核：它能够提高算术强度，并降低 LLM 预填充阶段的延迟。
最后讨论 KV cache 的内存需求，以及在使用 batching 提高推理阶段 GPU 利用率时，
降低这类内存需求的技术。

读者应当已经对 GPU 上 LLM 实现的基础知识有了清晰理解。完成本章后，还有非常大的
设计空间值得继续探索。

\section{练习}
\label{sec:ch20-exercises}

\begin{enumerate}[label=\arabic*.,leftmargin=2.5em]
  \item 为图 \ref{fig:20-9} 中的 FlashAttention 内核补充主机端调用代码。设置输入
  和输出数组的大小时，参考图 \ref{fig:20-8} 中所示矩阵的维度。
  \item 阅读 CUB 文档，了解 \code{WarpReduce} 命名空间中数据类型的细节，以及调用
  \code{Reduce} 函数时各参数的语义。特别注意归约操作所支持的各种算术函数。
  \item 编写图 \ref{fig:20-9} 第 21 行所调用的 \code{initialize()} device
  function，并解释为什么必须把 \code{O_i} 和 \code{D_i} 的元素初始化为
  \code{0.f}。还应说明为什么 \code{m_i} 需要初始化为负无穷。
  \item 证明图 \ref{fig:20-11} 的 device function 中嵌套循环对 $K$ 和 $V$ 的
  所有加载都是合并访问（coalesced）。
\end{enumerate}

\begin{chapterbibliography}{17}
  \bibitem{ch20-1}
  A. Vaswani, N. Shazeer, N. Parmar, J. Uszkoreit, L. Jones, A. N. Gomez,
  L. Kaiser, I. Polosukhin, \emph{Attention is all you need}, Proceedings of the
  31st International Conference on Neural Information Processing Systems, 2017,
  pp.~6000--6010.

  \bibitem{ch20-2}
  D. Ferrucci, E. Brown, J. Chu-Carroll, J. Fan, D. Gondek, A. A. Kalyanpur,
  A. Lally, J. W. Murdock, E. Nyberg, J. Prager, et al., \emph{Building Watson:
  an overview of the DeepQA project}, AI Magazine 31 (3), 2010, pp.~59--79.

  \bibitem{ch20-3}
  OpenAI, \emph{ChatGPT}, \url{https://chat.openai.com/}, 2025
  (accessed 25 January 2025).

  \bibitem{ch20-4}
  A. Gu, T. Dao, \emph{Mamba: linear-time sequence modeling with selective state spaces},
  arXiv preprint arXiv:2312.00752, 2023.

  \bibitem{ch20-5}
  M. Poli, S. Massaroli, E. Nguyen, D. Y. Fu, T. Dao, S. Baccus, Y. Bengio,
  S. Ermon, C. Ré, \emph{Hyena hierarchy: towards larger convolutional language models},
  Proceedings of the International Conference on Machine Learning, PMLR, 2023,
  pp.~28043--28078.

  \bibitem{ch20-6}
  T. Dao, D. Fu, S. Ermon, A. Rudra, C. Ré, \emph{FlashAttention: fast and memory-efficient
  exact attention with IO-awareness}, Advances in Neural Information Processing Systems
  35, 2022, pp.~16344--16359.

  \bibitem{ch20-7}
  M. Milakov, N. Gimelshein, \emph{Online normalizer calculation for softmax},
  arXiv preprint arXiv:1805.02867, 2018.

  \bibitem{ch20-8}
  T. Dao, \emph{FlashAttention-2: faster attention with better parallelism and work
  partitioning}, arXiv preprint arXiv:2307.08691, 2023.

  \bibitem{ch20-9}
  J. Shah, G. Bikshandi, Y. Zhang, V. Thakkar, P. Ramani, T. Dao,
  \emph{FlashAttention-3: fast and accurate attention with asynchrony and low-precision},
  arXiv preprint arXiv:2407.08608, 2024.

  \bibitem{ch20-10}
  Z. Zhang, D. Zhao, X. Miao, G. Oliaro, Q. Li, Y. Jiang, Z. Jia,
  \emph{Quantized side tuning: fast and memory-efficient tuning of quantized large language
  models}, arXiv preprint arXiv:2401.07159, 2024.

  \bibitem{ch20-11}
  J. Zhao, Z. Zhang, B. Chen, Z. Wang, A. Anandkumar, Y. Tian,
  \emph{Galore: memory-efficient LLM training by gradient low-rank projection},
  arXiv preprint arXiv:2403.03507, 2024.

  \bibitem{ch20-12}
  W. Kwon, Z. Li, S. Zhuang, Y. Sheng, L. Zheng, C. H. Yu, J. Gonzalez,
  H. Zhang, I. Stoica, \emph{Efficient memory management for large language model serving
  with PagedAttention}, Proceedings of the 29th Symposium on Operating Systems Principles,
  2023, pp.~611--626.

  \bibitem{ch20-13}
  T. Zadouri, H. Strauss, T. Dao, \emph{Hardware-efficient attention for fast decoding},
  arXiv preprint arXiv:2505.21487, 2025.

  \bibitem{ch20-14}
  Y. Leviathan, M. Kalman, Y. Matias, \emph{Fast inference from transformers via
  speculative decoding}, Proceedings of the International Conference on Machine Learning,
  PMLR, 2023, pp.~19274--19286.

  \bibitem{ch20-15}
  J. Ainslie, J. Lee-Thorp, M. de Jong, Y. Zemlyanskiy, F. Lebrón, S. Sanghai,
  \emph{GQA: training generalized multi-query transformer models from multi-head checkpoints},
  arXiv preprint arXiv:2305.13245, 2023.

  \bibitem{ch20-16}
  A. Liu, B. Feng, B. Wang, B. Wang, B. Liu, C. Zhao, C. Deng, C. Ruan, D. Dai,
  D. Guo, et al., \emph{DeepSeek-V2: a strong, economical, and efficient mixture-of-experts
  language model}, arXiv preprint arXiv:2405.04434, 2024.

  \bibitem{ch20-17}
  A. Q. Jiang, A. Sablayrolles, A. Roux, A. Mensch, B. Savary, C. Bamford,
  D. S. Chaplot, D. de la Casas, E. B. Hanna, B. Bressand, et al.,
  \emph{Mixtral of experts}, arXiv preprint arXiv:2401.04088, 2024.
\end{chapterbibliography}
